% GENERATED FILE. Edit canonical lessons, then run npm run generate:books.
% canonical-source-hash: fnv1a64:7c5eb6f6aa61781f
% canonical-lessons: UR-C435-khalij, UR-C435-langar, UR-C435-jal, UR-C435-sip, UR-C435-dubki

\chapter{Bays and Nets}
\label{ch:ur-bays-and-nets}
\hlchaptermodality{435}

\begin{chapteropening}
\textbf{By the end of this chapter:} \emph{I can describe a bay, its boats and what the fishermen bring in.}

Complete the last lesson of chapter 435: I can describe a bay, its boats and what the fishermen bring in.
\end{chapteropening}

\section[\texorpdfstring{\ur{خلیج} (khalīj)}{خلیج (khalīj)}]{\ur{خلیج} (khalīj) — a bay; a gulf}

\label{lesson:UR-C435-khalij}



\begin{quote}
Before the new one: say the Urdu for to toss (up), then the Urdu for to dip.
\end{quote}

\subsection*{You'll want to know: \ur{خلیج}}
\textbf{\ur{خلیج}} — \emph{khalīj} — \textquotedblleft{}a bay; a gulf\textquotedblright{}.

\begin{quote}
\ur{خلیج} \ur{کا} \ur{پانی} \ur{گہرا} \ur{اور} \ur{نیلا} \ur{ہے۔}
\end{quote}

\begin{quote}
\emph{khalīj kā pānī gahrā aur nīlā hai.}
\end{quote}

\begin{quote}
The water in the bay is deep and blue.
\end{quote}

Feminine.

\begin{cousinweb}
\textbf{\ur{خلیج}} (\emph{khalīj}) is usually traced, through Persian, to Arabic \emph{khalīj}, \textquotedblleft{}a bay, a gulf\textquotedblright{}.
\end{cousinweb}

\begin{grammarlens}[title={What you've built}]
One more word for getting around.
\end{grammarlens}

\subsection*{Your turn}
\begin{itemize}
\raggedright
  \item \emph{Say it:} \emph{khalīj}
  \item \emph{Say it:} \emph{khalīj}, once more
  \item \emph{Say it:} \emph{khalīj kā pānī gahrā aur nīlā hai.} — The water in the bay is deep and blue.
  \item \emph{Recall:} say the Urdu for to toss (up), then the Urdu for to dip, then say \emph{khalīj} again
\end{itemize}

\begin{tcolorbox}[breakable,colback=teal!4,colframe=teal!35!black,title={Before you move on}]
What does \ur{خلیج} mean? (\textquotedblleft{}A bay; a gulf\textquotedblright{}.) Say it once more.
\end{tcolorbox}
\section[\texorpdfstring{\ur{لنگر} (langar)}{لنگر (langar)}]{\ur{لنگر} (langar) — an anchor}

\label{lesson:UR-C435-langar}



\begin{quote}
Before the new one: say the Urdu for to dip, then the Urdu for a bay.
\end{quote}

\subsection*{You'll want to know: \ur{لنگر}}
\textbf{\ur{لنگر}} — \emph{langar} — \textquotedblleft{}an anchor\textquotedblright{}.

\begin{quote}
\ur{ملاحوں} \ur{نے} \ur{لنگر} \ur{اٹھایا۔}
\end{quote}

\begin{quote}
\emph{mallāhoṅ ne langar uṭhāyā.}
\end{quote}

\begin{quote}
The sailors weighed anchor.
\end{quote}

Masculine.

\begin{cousinweb}
\textbf{\ur{لنگر}} (\emph{langar}) is usually traced to Persian \emph{langar}, \textquotedblleft{}an anchor\textquotedblright{}; the same word names a free kitchen.
\end{cousinweb}

\begin{grammarlens}[title={What you've built}]
One more word for getting around.
\end{grammarlens}

\subsection*{Your turn}
\begin{itemize}
\raggedright
  \item \emph{Say it:} \emph{langar}
  \item \emph{Say it:} \emph{langar}, once more
  \item \emph{Say it:} \emph{mallāhoṅ ne langar uṭhāyā.} — The sailors weighed anchor.
  \item \emph{Recall:} say the Urdu for to dip, then the Urdu for a bay, then say \emph{langar} again
\end{itemize}

\begin{tcolorbox}[breakable,colback=teal!4,colframe=teal!35!black,title={Before you move on}]
What does \ur{لنگر} mean? (\textquotedblleft{}An anchor\textquotedblright{}.) Say it once more.
\end{tcolorbox}
\section[\texorpdfstring{\ur{جال} (jāl)}{جال (jāl)}]{\ur{جال} (jāl) — a net}

\label{lesson:UR-C435-jal}



\begin{quote}
Before the new one: say the Urdu for a bay, then the Urdu for an anchor.
\end{quote}

\subsection*{You'll want to know: \ur{جال}}
\textbf{\ur{جال}} — \emph{jāl} — \textquotedblleft{}a net\textquotedblright{}.

\begin{quote}
\ur{مچھیرے} \ur{نے} \ur{دریا} \ur{میں} \ur{جال} \ur{پھینکا۔}
\end{quote}

\begin{quote}
\emph{machhere ne daryā meṅ jāl pheṅkā.}
\end{quote}

\begin{quote}
The fisherman cast his net into the river.
\end{quote}

Masculine.

\begin{cousinweb}
\textbf{\ur{جال}} (\emph{jāl}) is inherited: CDIAL lists the Hindi cousin \emph{jāl}, \textquotedblleft{}a net for catching birds or fish\textquotedblright{}, under Sanskrit \emph{jāla} (CDIAL 5213).
\end{cousinweb}

\begin{grammarlens}[title={What you've built}]
One more word for getting around.
\end{grammarlens}

\subsection*{Your turn}
\begin{itemize}
\raggedright
  \item \emph{Say it:} \emph{jāl}
  \item \emph{Say it:} \emph{jāl}, once more
  \item \emph{Say it:} \emph{machhere ne daryā meṅ jāl pheṅkā.} — The fisherman cast his net into the river.
  \item \emph{Recall:} say the Urdu for a bay, then the Urdu for an anchor, then say \emph{jāl} again
\end{itemize}

\begin{tcolorbox}[breakable,colback=teal!4,colframe=teal!35!black,title={Before you move on}]
What does \ur{جال} mean? (\textquotedblleft{}A net\textquotedblright{}.) Say it once more.
\end{tcolorbox}
\section[\texorpdfstring{\ur{سیپ} (sīp)}{سیپ (sīp)}]{\ur{سیپ} (sīp) — a seashell}

\label{lesson:UR-C435-sip}



\begin{quote}
Before the new one: say the Urdu for an anchor, then the Urdu for a net.
\end{quote}

\subsection*{You'll want to know: \ur{سیپ}}
\textbf{\ur{سیپ}} — \emph{sīp} — \textquotedblleft{}a seashell\textquotedblright{}.

\begin{quote}
\ur{میں} \ur{نے} \ur{ساحل} \ur{سے} \ur{ایک} \ur{سیپ} \ur{اٹھائی۔}
\end{quote}

\begin{quote}
\emph{maiṅ ne sāhil se ek sīp uṭhāī.}
\end{quote}

\begin{quote}
I picked up a shell on the beach.
\end{quote}

Feminine.

\begin{cousinweb}
\textbf{\ur{سیپ}} (\emph{sīp}) is inherited: CDIAL lists the Hindi cousin \emph{sīp}, \textquotedblleft{}a bivalve shell\textquotedblright{}, under \emph{sippī}, a form it reconstructs (CDIAL 13417).
\end{cousinweb}

\begin{grammarlens}[title={What you've built}]
One more word for getting around.
\end{grammarlens}

\subsection*{Your turn}
\begin{itemize}
\raggedright
  \item \emph{Say it:} \emph{sīp}
  \item \emph{Say it:} \emph{sīp}, once more
  \item \emph{Say it:} \emph{maiṅ ne sāhil se ek sīp uṭhāī.} — I picked up a shell on the beach.
  \item \emph{Recall:} say the Urdu for an anchor, then the Urdu for a net, then say \emph{sīp} again
\end{itemize}

\begin{tcolorbox}[breakable,colback=teal!4,colframe=teal!35!black,title={Before you move on}]
What does \ur{سیپ} mean? (\textquotedblleft{}A seashell\textquotedblright{}.) Say it once more.
\end{tcolorbox}
\section[\texorpdfstring{\ur{ڈبکی} (ḍubkī)}{ڈبکی (ḍubkī)}]{\ur{ڈبکی} (ḍubkī) — a dip (under water)}

\label{lesson:UR-C435-dubki}



\begin{quote}
Before the new one: say the Urdu for a net, then the Urdu for a seashell.
\end{quote}

\subsection*{You'll want to know: \ur{ڈبکی}}
\textbf{\ur{ڈبکی}} — \emph{ḍubkī} — \textquotedblleft{}a dip (under water)\textquotedblright{}.

\begin{quote}
\ur{سارا} \ur{نے} \ur{جھیل} \ur{میں} \ur{ایک} \ur{ڈبکی} \ur{لی۔}
\end{quote}

\begin{quote}
\emph{sārā ne jhīl meṅ ek ḍubkī lī.}
\end{quote}

\begin{quote}
Sara ducked under the water in the lake.
\end{quote}

Feminine.

\begin{cousinweb}
\textbf{\ur{ڈبکی}} (\emph{ḍubkī}) is built on \emph{ḍub-}, \textquotedblleft{}sink\textquotedblright{}, from \emph{ḍubb}, a form CDIAL reconstructs (CDIAL 5561); \textbf{\ur{ڈبکی} \ur{لگانا}} is to take a dip.
\end{cousinweb}

\begin{grammarlens}[title={What you've built}]
One more word for getting around.
\end{grammarlens}

\subsection*{Your turn}
\begin{itemize}
\raggedright
  \item \emph{Say it:} \emph{ḍubkī}
  \item \emph{Say it:} \emph{ḍubkī}, once more
  \item \emph{Say it:} \emph{sārā ne jhīl meṅ ek ḍubkī lī.} — Sara ducked under the water in the lake.
  \item \emph{Recall:} say the Urdu for a net, then the Urdu for a seashell, then say \emph{ḍubkī} again
\end{itemize}

\begin{tcolorbox}[breakable,colback=teal!4,colframe=teal!35!black,title={Before you move on}]
What does \ur{ڈبکی} mean? (\textquotedblleft{}A dip (under water)\textquotedblright{}.) Say it once more.
\end{tcolorbox}
